\documentclass[a4paper,10pt]{article}

\usepackage[utf8]{inputenc}
\usepackage[T1]{fontenc}
\usepackage[portuguese]{babel}
\usepackage[margin=0.85in]{geometry}
\usepackage{titlesec}
\usepackage{enumitem}
\usepackage{hyperref}
\usepackage{xcolor}
\usepackage{indentfirst}

\titleformat{\section}{\color{blue}\large\bfseries}{}{0em}{}[\titlerule]
\pagenumbering{gobble}
\setlist[itemize]{leftmargin=1.5em, itemsep=0.15em, topsep=0.2em}

\begin{document}

\begin{center}
    {\LARGE \textbf{Pedro Henrique Klein}}\\
    \vspace{0.2cm}
    \href{mailto:pedrohenriquekleinphg@gmail.com}{pedrohenriquekleinphg@gmail.com} \textbar\
    (54) 99704-9008 \textbar\
    \href{https://www.linkedin.com/in/pedro-henrique-klein-a41122221/}{LinkedIn} \textbar\
    \href{https://github.com/Pedrinhonitz}{GitHub} \textbar\
    Chapecó, SC
\end{center}

\vspace{0.3cm}

\section*{Sobre}

Estudei programação desde os 15 anos, primeiro com Arduino e depois com Python, por conta própria. Aos 16 entrei no mercado como desenvolvedor e, até os 18, trabalhei em backend web com Django, Node.js, PHP, JavaScript e TypeScript.

Em 2021 ingressei em Ciência da Computação na Universidade Federal da Fronteira Sul (UFFS), em Chapecó, onde sigo estudando. No início de 2022 saí do backend e fui para engenharia de dados. Desde então construo pipelines, modelo dados e orquestro processos com Python, SQL, Airflow, GCP e AWS.

Essa mistura de desenvolvimento e dados me deixa transitar entre áreas técnicas e trabalhar bem com times diferentes.

\section*{Formação acadêmica}

\textbf{Universidade Federal da Fronteira Sul (UFFS)} \hfill 2021 a 2027 \\
Bacharelado em Ciência da Computação, campus Chapecó

\section*{Experiência profissional}

\textbf{BIP Brasil, projeto Pro Carbono (Bayer)} \hfill Set 2024 a Ago 2026

\textit{Engenheiro de Dados Sênior} \hfill Ago 2025 a Ago 2026
\begin{itemize}
    \item Implemento e mantenho pipelines com Apache Airflow, PySpark, AWS Glue e Astronomer Cosmos.
    \item Crio aplicações e APIs internas com Python, FastAPI e TypeScript para automatizar processos e integrar sistemas.
    \item Modelo dados com DBT, com qualidade, organização e rastreabilidade das transformações.
    \item Administro e otimizo ambientes Linux, Docker e Kubernetes para workloads de engenharia de dados.
    \item Escrevo e otimizo consultas em PostgreSQL, Amazon Athena e BigQuery para grandes volumes de dados.
    \item Estruturo esteiras de CI/CD, reviso código e defino práticas de desenvolvimento.
    \item Contribuo com o Apache Airflow, com correções e melhorias para a comunidade.
\end{itemize}

\textit{Engenheiro de Dados Pleno} \hfill Set 2024 a Ago 2025
\begin{itemize}
    \item Implementei e mantive pipelines com Apache Airflow e PySpark.
    \item Construí soluções de processamento e armazenamento em AWS e Google Cloud Platform.
    \item Modelei dados com DBT e Elementary.
    \item Implantei CI/CD com GitHub Actions, Jenkins e Terraform, além de padrões de código e revisão de pull requests.
\end{itemize}

\textbf{Descomplica} \hfill Jan 2022 a Set 2024 \\
\textit{Engenheiro de Dados Pleno}
\begin{itemize}
    \item Orquestro workflows no Apache Airflow para executar e monitorar pipelines de dados.
    \item Empacoto cargas com Docker e orquestro workloads no Kubernetes.
    \item Administro e otimizo ambientes Linux.
    \item Construo soluções de data warehouse com BigQuery, Amazon Athena e Snowflake.
    \item Monto dashboards e relatórios no Metabase.
    \item Automatizo ETL com Python e Shell.
    \item Escrevo e otimizo SQL para grandes volumes de dados.
\end{itemize}

\textbf{IXC-Soft} \hfill Nov 2021 a Jan 2022 \\
\textit{Estagiário em Engenharia de Software}
\begin{itemize}
    \item Desenvolvi backend em PHP.
    \item Montei interfaces com Twig.
    \item Empacotei ambientes de desenvolvimento com Docker.
    \item Administrei e otimizei bancos MariaDB.
\end{itemize}

\section*{Open Source}

\textbf{\href{https://github.com/apache/airflow/pull/66431}{Apache Airflow (Apache Software Foundation)}} \\
Pull request aceito no repositório oficial.
\begin{itemize}
    \item Corrigi o scheduler no tratamento de \textit{deferred tasks}, removendo uma condição de corrida com eventos obsoletos enviados pelo executor.
    \item Trabalhei com os mantenedores na revisão técnica até a mudança ser aprovada e integrada ao código principal.
\end{itemize}

\section*{Habilidades técnicas}

\begin{itemize}
    \item \textbf{Linguagens:} Python, SQL, TypeScript, JavaScript, Bash, Go, Java, C, C\texttt{++}, Haskell e Ruby.
    \item \textbf{Engenharia de dados:} Apache Airflow, Astronomer Cosmos, DBT, Elementary, PySpark, AWS Glue, BigQuery, Snowflake, Amazon Athena, Luigi e Dagster.
    \item \textbf{DevOps e infraestrutura:} Docker, Kubernetes, Terraform, GitHub Actions, Jenkins e Linux.
    \item \textbf{Bancos de dados:} PostgreSQL, SQL Server, Oracle e Neo4j.
    \item \textbf{Cloud:} Amazon Web Services (AWS) e Google Cloud Platform (GCP).
\end{itemize}

\section*{Idiomas}

\begin{itemize}
    \item Português: nativo
    \item Inglês: A2
    \item Espanhol: B1
\end{itemize}

\end{document}
